% Generated draft: requires mathematical and source review.
\documentclass[11pt]{article}
\usepackage[margin=25mm]{geometry}
\usepackage{fontspec}
\setmainfont{Latin Modern Roman}
\usepackage{amsmath,amssymb,mathtools,unicode-math}
\setmathfont{Latin Modern Math}
\usepackage{enumitem,needspace,xurl,hyperref}
\setlength{\parindent}{0pt}
\setlength{\parskip}{5pt}
\setlength{\emergencystretch}{3em}
\newcommand{\N}{\mathbb{N}}
\newcommand{\Z}{\mathbb{Z}}
\newcommand{\Q}{\mathbb{Q}}
\newcommand{\R}{\mathbb{R}}
\newcommand{\C}{\mathbb{C}}
\newcommand{\F}{\mathbb{F}}
\newcommand{\A}{\mathbb{A}}
\newcommand{\PP}{\mathbb{P}}
\newcommand{\E}{\mathbb{E}}
\newcommand{\eps}{\varepsilon}
\newcommand{\dd}{\,\mathrm{d}}
\begin{document}
\section*{UnsolvedMath ID: 30004431}
\textbf{Source problem code:} catalogue entry
\subsection{Definitions and mathematical statement}
Let \(f\) denote the potential and \(\mu\) an equilibrium state referred to in the supplied record. The underlying dynamical system, the class of potentials, and the transfer operator are not specified. Consequently, the following records the intended questions with their unresolved parameters; it is not a fully specified equivalence conjecture.

In a dynamical setting \((X,T)\) where the variational definition applies, write \(\mathcal M_T(X)\) for the admissible \(T\)-invariant probability measures. An equilibrium state for \(f\) is a measure \(\mu\in\mathcal M_T(X)\) satisfying
\[
 h_\mu(T)+\int_X f\,\dd\mu
 =\sup_{\nu\in\mathcal M_T(X)}\left(h_\nu(T)+\int_X f\,\dd\nu\right),
\]
where \(h_\nu(T)\) is measure-theoretic entropy, provided these expressions are defined. This is a conditional explanation of the terminology, not a choice of the missing system or admissibility hypotheses.

Writing \(\mathcal L_f\) for the intended transfer operator, a strictly positive continuous eigenfunction means a function \(h\in C(X,\R)\) and an eigenvalue \(\lambda\) such that
\[
 h(x)>0\quad(x\in X),\qquad \mathcal L_f h=\lambda h.
\]
The topology on \(X\), the definition and domain of \(\mathcal L_f\), and any required normalization or spectral condition on \(\lambda\) remain unspecified.

A \(g\)-measure ordinarily refers to a stationary measure whose one-sided conditional probabilities are prescribed by a normalized function \(g\). Making this definition precise requires a symbol space, a conditioning direction, a normalization rule, and a specified regularity class for \(g\). The supplied record does not provide these choices or say whether \(g\) is fixed, is determined by \(f\), or is existentially quantified. Likewise, the term Gibbs measure requires a specified potential or interaction and a specified Gibbs condition; these are not supplied.

\textbf{Questions recorded by the source.} For an equilibrium state \(\mu\) of a potential \(f\), is
\[
 \bigl[\mathcal L_f\text{ has a strictly positive continuous eigenfunction}\bigr]
 \quad\Longleftrightarrow\quad
 \bigl[\mu\text{ is a }g\text{-measure}\bigr]?
\]
More generally, do the classes of Gibbs measures and \(g\)-measures coincide in the intended setting? The first question also leaves unresolved whether the equivalence concerns each equilibrium state, existence of a suitable equilibrium state, or an equilibrium state selected by additional transfer-operator data. No universal quantification over systems, potentials, or choices of \(g\) can be fixed from this record alone.

The supplied dated curation describes the problem as partially solved: it reports a positive continuous transfer-operator eigenfunction as yielding a \(g\)-measure, reports that converse and Gibbs--\(g\)-measure equivalences require additional regularity and have counterexamples in general, and describes the unrestricted equivalence as not solved. These are supplied assessments, not independently verified results here; their precise hypotheses are absent. In particular, existence of a resulting \(g\)-measure must not silently be identified with the assertion that the particular equilibrium state \(\mu\) is that measure.
\subsection{Short English statement}
The source asks whether a positive continuous transfer-operator eigenfunction characterizes when an equilibrium state is a \(g\)-measure, and whether Gibbs measures and \(g\)-measures coincide. The dynamical setting, regularity assumptions, Gibbs convention, and relation between \(f\), \(g\), and the chosen measure are missing. This statement-only draft preserves those gaps and the supplied partially solved classification for human mathematical review.
\subsection{Sources}
\begin{itemize}
\item Local registry: \texttt{data/unsolvedmath/problems.json}, with \texttt{catalog/problems.json} as the supplementary registry specified by the project; supplied record with permanent UnsolvedMath ID \(30004431\), title ``Equivalence of Gibbs Measures and \(g\)-Measures,'' exact source problem code \texttt{catalogue entry}, in the Oberwolfach Reports Open Problems collection. Registry files were not inspected for this draft.
\item Supplied report citation: \emph{Mini-Workshop: One-sided and Two-sided Stochastic Descriptions} (2021), Oberwolfach Reports, pp.\ 601--637, DOI \texttt{10.4171/owr/2020/11}. \url{https://doi.org/10.4171/owr/2020/11}. The record assigns proposed year 2020 and describes the report as documenting workshop open-problem discussions; no individual proposer is identified.
\item Supplied literature reference: Walters, \emph{Regularity Conditions and Bernoulli Properties of Equilibrium States and g-Measures}. \url{https://doi.org/10.1112/S0024610704006076}.
\item Supplied literature reference: Bissacot, Endo and Van Enter, \emph{Regular g-measures are not always Gibbsian}. \url{https://arxiv.org/abs/1106.4188}.
\item Supplied literature reference: OWR-2020-11. \url{https://ems.press/content/serial-article-files/46847}.
\item Literature assessment date: \texttt{2026-08-22}. The partially solved status and accompanying claims are supplied dated curation derived from the supplied OWR curation, not independently verified here. The record itself requests independent verification of time-sensitive status and cited claims.
\item The supplied statement status is \texttt{corrected\_verified}, accompanied by a claim of manual semantic, typographic, and LaTeX review. This is source metadata, not independent mathematical verification of the formulation or its completeness.
\end{itemize}
\end{document}
